\section{Water Orientation}
The orientation of water within the aqueous/organic interface of the system was defined using the angles formed by molecular axes and the fixed long-axis of the system (perpendicular to the interfacial plane). The two molecular axes used are the water bisector, pointing from the hydrogen side to the oxygen, and the vector normal to the plane formed by the three atoms of the water molecule. Figure \ref{fig:water-angles} illustrates the angle profiles of interest relative to the system reference axis (i.e. Z-axis) as the molecule takes on various orientations.

% find references of the orientation of water near salt-interfaces for the above paragraph

\begin{figure}[h!]
\begin{center}
	\includegraphics[scale=0.26]{images/h2o-2dhistograms.png}
	\caption{Orientation profiles of interfacial water molecules. Two profiles were developed for characterizing interfacial water orientation based on the bisector and molecular normal vector angles formed with the system references axis (see Figure \ref{fig:water-angles}). Each profile shows the water population for each cosine of the respective angle as a function of distance to the Gibbs dividing surface (located at 0.0). The profiles in the left column are of the bisector vector, and those on the right are of the molecular normal vector. Neat-H$_2$O, NaCl, NaNO$_3$, and Na$_2$SO$_4$ system profiles are ordered from the top to bottom row, respectively.}
	\label{fig:system-orientations}
\end{center}
\end{figure}

Previous studies have provided a detailed overview of water orientation at the interface with both air and organic phases.\cite{McFearin2009,Hore2008,Fan2009,Wick2006c,Wick2008a} The topmost water layers are highly ordered because of their contact with the organic phase, and it has been suggested that ordering of both the organic and water molecules would lead to a field across the boundary of the interface.\cite{McFearin2009,Hore2008} This can influence charged species, and the ordering and orientation of the H-bond network. Our recent experimental SFG results suggest that the accumulation of charged ions leads to a field-screening that affects the orientation of waters in the topmost layers. This is complemented by the results of the current study.

\newcommand{\degree}{\ensuremath{^\circ}}

\begin{figure}[h!]
\begin{center}
	\includegraphics[scale=0.26]{images/h2o-angle-blurs.png}
	\caption{Depictions of possible orientations for different ranges of the bisector-reference angle $\theta$. A range of $0<\cos(\theta)<1$ ($0\degree<\theta<90\degree$ - left figure) suggests that one of the OH-bonds will always point into the aqueous phase, and the other will straddle the plane of the interface. In the range $-0.5<\cos(\theta)<0.5$ ($60\degree<\theta<120\degree$ - center figure) one OH-bond will always point out of the aqueous phase, and the other will always remain pointing in towards the water bulk. The molecular normal angle, $\phi$, has a narrow range for interfacial waters in all of the systems spanning approx. $0.7<\cos(\theta)<1$ (right). The range of $\phi$ keeps the topmost waters lying mostly flat on the interface, but with a narrow range of twist with OH-bonds pointing both into, and out of the interface.}
	\label{fig:angle-blurs}
\end{center}
\end{figure}

The far-left dark regions of the plots in figure \ref{fig:water-angles} show the CCl$_4$ bulk near to the interface where few or no waters are found. The gibbs dividing surface is locate at 0.0, to the far right, in the water bulk, the flat profile represents an isotropic orientation of the waters, showing no particular directional or orientational preference. The regions of interest lie around the gibbs dividing surface within the interface. In the top plot of the neat-water reference system, the bisector orientation profile relative to the other systems is a bit anomolous. The bisector exhibits an affinity to point such that $0<\cos(\theta)<1$. The orientation of the OH-bonds for such an angle range is illustrated in figure \ref{fig:angle-blurs}. In the case of the neat-water reference system, the ``free-OH'' bond is generally pointing out of the aqueous phase, but often points within the plane of the interface. At $3\AA$ into the neat-water surface, there appears a peak spanning approx. $0.7<\cos(\theta)<0.9$ that stretches out to just past $5\AA$, at which point the profile becomes flat across the values of $\theta$.

The molecular plane normal in the right column of figure \ref{fig:water-angles} is clearly centered about the flatter orientations with $0.7<\cos(\theta)<1$. This profile suggests that the water lies flat on the interface with a twist of up to 45\degree.  Molecular normal orientation profiles are very similar across the systems investigated, with each profile showing an affinity for the flat orientation starting just before the gibbs dividing surface, and extending as deep as $7\AA$ in the water and NaCl systems, and between $3-4\AA$ in the NaNO$_3$ and Na$_2$SO$_4$ polyatomic salt systems.

Bisector profiles for the salt systems are shifted from that of the neat-water. Added salt ions appear to rotate the interfacial water orientation such that the OH-bond that would straddle the interface in the neat-water system points out of the interface with a greater angle. The inner OH-bond also points into the aqueous phase, but with a more narrow range of orientations. The profiles of the salt systems all center about the $\cos(\theta)=0$ region that indicates a straddling water molecule. All three profiles are also more heavily weighted to $\cos(\theta)<0$, whereas the neat-water system followed the opposite trend. The orientation effect is shallowest in the NaCl system, and is felt deepest in the NaNO$_3$. Although the Na$_2$SO$_4$ profile becomes isotropic at nearly the same depth as the NaCl, there is a clear difference in a shift around $5\AA$ of orientation away from the $\cos(\theta)>0$ that lasts until almost $8\AA$. The deeper waters are more clearly oriented on the bottom half of the profile in the Na$_2$SO$_4$ system.

The sharpest transition in the bisector profiles is that of the Na$_2$SO$_4$ system. This can be attributed to the proximity of the anion and cation density peaks in the water bulk, which is shortest in the system. Following the trend is the broadest bisector profile of NaNO$_3$ that extends past $10\AA$ past the interface, and corresponds to the greatest anion-cation density peak distance. NaCl follows the trend as the bisector orientation profile extends to the same depth as the Na$_2$SO$_4$ system, but the transition to the isotropic region is not nearly as sharp, as the electric double-layer set up by the anion and cation densities is not as strong. This is both because of the proximity of the ions, but also because of the higher charge of the divalent SO$_4^{2-}$ anion compared to the monovalent Cl$^-$. The field created by the anion-cation pairing within the interface affects both the depth to which waters are oriented, and also the width of the transition to isotropic orientation.

